\begin{abstract}
Two lines of work try to make language models more human-like: architectures that limit memory, as human working memory is limited, and training on child-scale input, as in the BabyLM Challenge. Do these constraints bring models closer to humans? We compare the representational geometry of 118 language models, crossing a memory constraint with a data constraint, with that of 16 reading-time measures and eight fMRI, EEG, MEG and intracranial datasets. At matched size and perplexity, both constraints are associated with higher similarity to reading behavior but not to brain activity, where standard large language models remain the most similar. We then ask where the behavioral advantage comes from. Syntactically demanding words give only a weak test: once length and frequency are controlled, they slow readers in a single corpus, and only slightly, and there the constraint advantage is not reliably larger than on easy words. The advantage is concentrated on rare words, where reading times hold most of their model-explainable signal, although not in every reading measure. On the neural side, constrained models gain no advantage on the domain-general multiple-demand network, which supports working memory but which no model predicts well. Human-inspired constraints are thus associated with more human-like word difficulty, but not with more brain-like representations.
\end{abstract}
